% sections/anomaly_injection.tex — Anomaly Injection Framework
\section{Anomaly Injection Framework}
\label{sec:anomaly_injection}

A primary use case for \system{} is generating labeled datasets for fraud
detection and anomaly analytics.  This section describes the multi-dimensional
anomaly injection framework.

% ─────────────────────────────────────────────────────────────────────
\subsection{Anomaly Taxonomy}
\label{sec:anom_taxonomy}

\system{} defines 33 anomaly types organized into five categories:

\begin{table}[t]
\centering
\caption{Anomaly taxonomy with representative types per category.}
\label{tab:anomaly_types}
\small
\begin{tabular}{p{2.4cm}cp{7.5cm}}
\toprule
\textbf{Category} & \textbf{Count} & \textbf{Representative Types} \\
\midrule
Fraud & 10 & FictitiousEntry, RoundDollarManipulation,
  JustBelowThreshold, SelfApproval, DuplicatePayment,
  FictitiousVendor, RevenueManipulation \\
Error & 8 & DuplicateEntry, ReversedAmount, WrongPeriod,
  WrongAccount, MissingReference \\
Process & 5 & LatePosting, SkippedApproval,
  ThresholdManipulation \\
Statistical & 5 & UnusualAmount, TrendBreak,
  BenfordViolation \\
Relational & 5 & CircularTransaction,
  DormantAccountActivity \\
\bottomrule
\end{tabular}
\end{table}

Each anomaly type is parameterized by:
\begin{itemize}
  \item \textbf{Weight:} Relative injection probability (e.g., 2.5 for
    JustBelowThreshold vs.\ 0.5 for FictitiousVendor).
  \item \textbf{Amount range:} Domain-specific bounds (e.g.,
    \$10K--\$500K for FictitiousEntry, \$100K--\$5M for
    RevenueManipulation).
  \item \textbf{Conditionality:} Some types require specific preconditions
    (e.g., SelfApproval requires matching approver and preparer).
\end{itemize}

% ─────────────────────────────────────────────────────────────────────
\subsection{Multi-Stage Fraud Schemes}
\label{sec:anom_schemes}

Simple anomalies are injected as independent point events.  However, real
fraud often evolves through stages over extended time horizons.  \system{}
implements three multi-stage \emph{fraud schemes}:

\paragraph{1. Vendor Kickback Scheme.}
A buyer routes payments to a shell company they control.  The scheme
progresses through stages: (a)~establish shell vendor, (b)~route initial
small invoices, (c)~gradually increase amounts, (d)~introduce collusion
patterns across transactions.  Stage transitions are governed by
probability-weighted state machines.

\paragraph{2. Gradual Embezzlement Scheme.}
Small, individually undetectable amounts are siphoned over time.  The scheme
models: (a)~initial testing with minimal amounts, (b)~establishment of a
regular extraction pattern, (c)~acceleration in amount and frequency,
(d)~period-end concentration where oversight is reduced.

\paragraph{3. Revenue Manipulation Scheme.}
Fictitious customers are created and revenue is prematurely recognized.
Stages include: (a)~customer creation, (b)~fictitious order entry,
(c)~premature revenue recognition, (d)~channel stuffing patterns with
period-end concentration.

The \texttt{SchemeAdvancer} tracks each scheme's state, selecting
context-aware actions at each stage and maintaining cumulative impact
metadata.  Each transaction within a scheme receives a
\texttt{MultiStageAnomalyLabel} that links it to the scheme instance,
stage, and causal chain.

% ─────────────────────────────────────────────────────────────────────
\subsection{Anomaly Correlation Patterns}
\label{sec:anom_correlation}

Anomalies in real data are not independently distributed.  \system{} models
three correlation mechanisms:

\paragraph{Temporal Clustering.}
Anomalies concentrate in time windows (e.g., period-end, after management
changes).  The \texttt{TemporalClusterGenerator} defines windows with
configurable intensity multipliers:

\begin{equation}
P(\text{anomaly at } t) = p_0 \cdot
  \prod_{w \in \mathcal{W}} \left(1 + (m_w - 1) \cdot
  \mathbb{1}_{t \in w}\right)
\end{equation}

where $p_0$ is the base anomaly rate and $m_w$ is the multiplier for window $w$.

\paragraph{Cascade Propagation.}
One anomaly triggers follow-on anomalies in dependent processes.  A
\texttt{CascadeStep} defines the trigger condition, dependent anomaly type,
propagation probability, and expected delay.  For example, a three-way
match variance (trigger) may cascade to a payment delay (dependent) with
probability 0.7 and a 5-day lag.

\paragraph{Co-Occurrence.}
Multiple anomaly types appear on the same transaction.  The
\texttt{CoOccurrencePattern} defines allowed combinations and prevents
unrealistic pairings (e.g., a FictitiousEntry cannot simultaneously be a
DuplicateEntry).

% ─────────────────────────────────────────────────────────────────────
\subsection{Difficulty, Severity, and Confidence Scoring}
\label{sec:anom_scoring}

Each injected anomaly receives three scores that characterize its
detectability:

\paragraph{Difficulty.}
A composite score $d \in [0, 1]$ (0 = easy to detect, 1 = hard) computed
from five factor groups:

\begin{equation}
d = \sum_{f \in \mathcal{F}} w_f \cdot s_f, \quad
\mathcal{F} = \{\text{amount}, \text{blending}, \text{collusion},
  \text{concealment}, \text{temporal}\}
\label{eq:difficulty}
\end{equation}

\begin{itemize}
  \item \textbf{Amount factors:} Larger amounts are easier to detect
    ($s_{\text{amount}}$ decreases with magnitude).
  \item \textbf{Blending factors:} How well the anomaly blends with normal
    data ($s_{\text{blending}} \in [0,1]$).
  \item \textbf{Collusion factors:} Multi-party schemes are harder to detect.
  \item \textbf{Concealment factors:} Obfuscation techniques (round-tripping,
    split transactions) increase difficulty.
  \item \textbf{Temporal factors:} Night/weekend postings, period-end timing
    reduce oversight.
\end{itemize}

\paragraph{Severity.}
A five-component model assessing organizational impact:

\begin{equation}
\text{severity} = 0.25 \cdot s_{\text{type}} + 0.30 \cdot s_{\text{monetary}}
  + 0.20 \cdot s_{\text{frequency}} + 0.15 \cdot s_{\text{scope}}
  + 0.10 \cdot s_{\text{timing}}
\label{eq:severity}
\end{equation}

where $s_{\text{type}}$ is the inherent risk per anomaly category,
$s_{\text{monetary}}$ normalizes monetary impact against materiality,
$s_{\text{frequency}}$ escalates for repeated patterns,
$s_{\text{scope}}$ captures the number of affected entities, and
$s_{\text{timing}}$ penalizes period-end and audit-period timing.

\paragraph{Confidence.}
Detection confidence $c \in [0,1]$ estimates the strength of evidence
available to an auditor or ML model.  It combines pattern-match confidence
(statistical correlation with known anomaly signatures), rule-violation
confidence (policy breaches), behavioral-deviation confidence (distance from
baseline), and contextual modifiers (co-occurring anomalies, scheme
membership).

% ─────────────────────────────────────────────────────────────────────
\subsection{Ground-Truth Labels}
\label{sec:anom_labels}

Every anomaly produces a structured label record containing:

\begin{itemize}
  \item Anomaly type, category, and subcategory.
  \item Affected transaction IDs and account codes.
  \item Difficulty, severity, and confidence scores.
  \item Scheme identifier and stage (for multi-stage schemes).
  \item Contributing factors with individual impact and reliability scores.
  \item Temporal context (posting timestamp, detection window).
\end{itemize}

Labels are exported as separate CSV/JSON files (\texttt{anomaly\_labels},
\texttt{fraud\_labels}, \texttt{quality\_labels}), enabling supervised,
semi-supervised, and unsupervised evaluation protocols.

% ─────────────────────────────────────────────────────────────────────
\subsection{Banking-Specific Anomalies: AML Typologies}
\label{sec:anom_aml}

The banking module extends the anomaly framework with five Anti-Money
Laundering (AML) typologies:

\begin{enumerate}
  \item \textbf{Structuring (Smurfing):} Multiple deposits below the
    reporting threshold (\$10{,}000 in the US), distributed across accounts
    and time windows.
  \item \textbf{Funnel Accounts:} Rapid aggregation of funds from multiple
    sources into a single intermediary account.
  \item \textbf{Layering:} Complex chains of wire transfers across
    jurisdictions to obscure the origin of funds.
  \item \textbf{Mule Networks:} Coordinated account networks receiving and
    forwarding stolen funds.
  \item \textbf{Round Tripping:} Funds exit and re-enter the same entity
    disguised as legitimate commerce, with 30--60 day delays.
\end{enumerate}

Each typology generates persona-based transaction patterns with risk-tier
profiles (Low, Medium, High, Very High) and sophistication levels (Basic,
Intermediate, Advanced).  Ground-truth AML labels are exported alongside
transaction data.

% ─────────────────────────────────────────────────────────────────────
\subsection{Fraud Scenario Packs}
\label{sec:anom_fraud_packs}

\system{} provides five pre-built fraud scenario packs as YAML configurations
that can be deep-merged into any generation config:

\begin{itemize}
  \item \texttt{revenue\_fraud}: Revenue manipulation patterns including
    channel stuffing, bill-and-hold, and premature recognition.
  \item \texttt{payroll\_ghost}: Ghost employee schemes with fictitious
    payroll entries.
  \item \texttt{vendor\_kickback}: Vendor collusion with inflated invoices
    and split payments.
  \item \texttt{management\_override}: Override of internal controls by
    senior personnel.
  \item \texttt{comprehensive}: Combined pack with elements from all four
    categories.
\end{itemize}

These packs configure anomaly rates, affected account ranges, temporal
distribution, and multi-stage progression.  They are applied via
\texttt{apply\_fraud\_packs()} which performs a deep-merge, and can be
overridden with \texttt{-{}-fraud-rate} for rate adjustment.
